% TOP_PROBLEM: {"id": 3317, "title": "3-Colourability of Arrangements of Great Circles", "category_id": 3, "category": {"id": 3, "name": "graph_theory", "display_name": "Graph Theory", "description": "Problems involving graphs, networks, and their properties.", "slug": "graph-theory", "order_index": 3, "created_at": "2026-07-31T15:26:25.671Z"}, "status": "open", "problem_number": "OPG-34915", "set_id": 12}
% FIRST_POSTED: 2026-10-01
% ENTRY_KIND: statement_only
% UnsolvedMath ID 3317; source catalogue code OPG-34915
% !TeX program = lualatex
% Definition and sources only; this file is not an LLM solution attempt.
\documentclass[11pt,a4paper]{article}
\usepackage[margin=25mm]{geometry}
\usepackage{fontspec}
\setmainfont{lmroman10-regular.otf}[BoldFont=lmroman10-bold.otf,
 ItalicFont=lmroman10-italic.otf,BoldItalicFont=lmroman10-bolditalic.otf]
\usepackage{amsmath,amssymb,mathtools}
\usepackage{unicode-math}
\setmathfont{latinmodern-math.otf}
\AtBeginDocument{\let\setminus\smallsetminus}
\usepackage{xurl,hyperref}
\hypersetup{colorlinks=true,urlcolor=blue}
\setcounter{secnumdepth}{0}
\setlength{\parindent}{0pt}
\setlength{\parskip}{5pt}
\setlength{\emergencystretch}{3em}
\begin{document}
\section*{3-Colourability of Arrangements of Great Circles}\label{3317}
{\small UnsolvedMath ID 3317 · OPG-34915 · Graph Theory · OpenGarden}
\subsection{Definitions and mathematical statement}
Let $S^2=\{x\in\mathbb R^3:\|x\|=1\}$. A great circle is the
intersection of $S^2$ with a two-dimensional linear subspace of
$\mathbb R^3$. Fix a finite collection $\mathcal C$ of $m\ge2$
distinct great circles in general position in the sense that no
point of the sphere belongs to three of them. Each pair meets at
two antipodal points.

The arrangement graph has these intersection points as vertices.
For each circle, join two consecutive intersection points along
that circle by the intervening circular arc, provided that arc has
no other intersection point in its interior. These arcs are the
edges of the embedded graph. The graph is 4-regular when parallel
edges are counted; in the two-circle case parallel edges do occur
and impose the same colouring constraint as a single edge.

The conjecture asks whether this arrangement graph always has a
proper 3-colouring: a map from its vertices to $\{1,2,3\}$ taking
different values at the endpoints of every arrangement edge.
No relation is required between colours at antipodal vertices.
For fewer than two circles there are no intersection vertices,
so the colouring requirement is vacuous.

The restriction to great circles is a geometric hypothesis, not
merely that the graph is planar and 4-regular. Likewise, the
no-triple-intersection assumption specifies which arrangement graph
is intended. The target is ordinary vertex colouring from one
three-colour palette; the stronger list-colouring question mentioned
in the source is separate.
\subsection{Short English statement}
Is the intersection graph drawn along any general-position arrangement of great circles properly vertex-colourable with three colours?
\subsection{Sources}
\begin{itemize}
\item[S1] ULAM.AI and the UnsolvedMath Contributors, supplied problems.json, record 3317 (OPG-34915), OpenGarden collection. Catalogue identity and source transcription; CC BY 4.0.
\url{https://huggingface.co/datasets/ulamai/UnsolvedMath}.
\item[S2] Open Problem Garden, 3-Colourability of Arrangements of Great Circles. Original problem and discussion; the mathematical formulation below is expanded from the supplied source record.
\url{http://www.openproblemgarden.org/op/3_colourability_of_arrangements_of_great_circles}.
\end{itemize}
\end{document}
